\documentclass[11pt]{article}
\usepackage[margin=2.6cm]{geometry}
\usepackage{amsmath,booktabs,siunitx}
\usepackage{fancyhdr}
\pagestyle{fancy}
\fancyhf{}
\lhead{Simple lab report}
\rhead{\thepage}
\title{Measuring the Cooling of Hot Water}
\author{Your Name}
\date{\today}
\begin{document}
\maketitle
\section{Question}
How does the temperature of a cup of hot water change during the first ten
minutes after pouring?

\section{Prediction}
Newton's cooling model suggests that the rate of cooling is approximately
proportional to the difference between the water and room temperatures:
\[
  \frac{dT}{dt}=-k(T-T_\mathrm{room}).
\]

\section{Method}
Pour \SI{200}{mL} of hot water into an uncovered cup. Measure the temperature
with the same probe every two minutes, keeping the probe depth and cup position
constant. Record room temperature at the start and end.

\section{Results}
\begin{table}[ht]
\centering
\caption{Temperature observations.}
\begin{tabular}{@{}S[table-format=2.0] S[table-format=2.1] S[table-format=2.1]@{}}
\toprule
{$t$ (min)} & {$T$ (\si{\celsius})} & {$T-T_\mathrm{room}$ (\si{\celsius})}\\
\midrule
0 & 86.0 & 64.0\\
2 & 77.4 & 55.4\\
4 & 70.3 & 48.3\\
6 & 64.2 & 42.2\\
8 & 59.0 & 37.0\\
10 & 54.6 & 32.6\\
\bottomrule
\end{tabular}
\end{table}

\section{Interpretation}
Plot $\ln(T-T_\mathrm{room})$ against time, estimate the gradient, and compare
the fitted model with the observations. Mention convection, evaporation, probe
response time, and any reading that was taken late.

\section{Conclusion}
The fitted cooling constant is $k=\rule{0.25\linewidth}{0.4pt}$ per minute.
The prediction was \rule{0.2\linewidth}{0.4pt}; the strongest evidence is
\rule{0.4\linewidth}{0.4pt}.
\end{document}
